\section{Subpreschemes and immersion morphisms}
\label{section:I.4}

\subsection{Subpreschemes.}
\label{subsection:I.4.1}

\begin{env}[4.1.1]
\label{I.4.1.1}
Since the notion of quasi-coherent sheaf
(\hyperref[0.5.1.3]{0, 5.1.3}) is local, a
quasi-coherent $\mathscr{O}_X$-module $\mathscr{F}$ on a prescheme $X$ can
be defined by the condition that, for every affine open set $V$ of $X$,
$\mathscr{F}|V$ is isomorphic to a sheaf associated with a
$\Gamma(V,\mathscr{O}_X)$-module (\hyperref[I.1.4.1]{1.4.1}). It is clear
that, on a prescheme $X$, the structure sheaf $\mathscr{O}_X$ is
quasi-coherent, and that kernels, cokernels and images of homomorphisms of
quasi-coherent $\mathscr{O}_X$-modules, and inductive limits and direct sums
of quasi-coherent $\mathscr{O}_X$-modules, are quasi-coherent
(\hyperref[I.1.3.7]{1.3.7} and \hyperref[I.1.3.9]{1.3.9}).
\end{env}

\begin{proposition}[4.1.2]
\label{I.4.1.2}
Let $X$ be a prescheme, and $\mathscr{I}$ a quasi-coherent sheaf of ideals
in $\mathscr{O}_X$. Then the support $Y$ of the sheaf
$\mathscr{O}_X/\mathscr{I}$ is closed, and if we denote by
$\mathscr{O}_Y$ the restriction of $\mathscr{O}_X/\mathscr{I}$ to $Y$,
$(Y,\mathscr{O}_Y)$ is a prescheme.
\end{proposition}

\oldpage[I]{120}
It obviously suffices (\hyperref[I.2.1.3]{2.1.3}) to consider the case where
$X$ is an affine scheme, and to show that in this case $Y$ is closed in
$X$ and is an \emph{affine scheme}. Indeed, if
$X=\operatorname{Spec}(A)$, we have $\mathscr{O}_X=\widetilde{A}$ and
$\mathscr{I}=\widetilde{\mathfrak{I}}$, where $\mathfrak{I}$ is an ideal of
$A$ (\hyperref[I.1.4.1]{1.4.1}); $Y$ is then equal to the closed subset
$V(\mathfrak{I})$ of $X$ and is identified with the prime spectrum of the ring
$B=A/\mathfrak{I}$ (\hyperref[I.1.1.11]{1.1.11}); moreover, if $\varphi$
is the canonical homomorphism $A\to B=A/\mathfrak{I}$, the direct image
${}^a\varphi_*(\widetilde{B})$ is canonically identified with the sheaf
$\widetilde{A}/\widetilde{\mathfrak{I}}=\mathscr{O}_X/\mathscr{I}$
(\hyperref[I.1.6.3]{1.6.3} and \hyperref[I.1.3.9]{1.3.9}), which completes
the proof.

We shall say that $(Y,\mathscr{O}_Y)$ is the \emph{sub-prescheme} of
$(X,\mathscr{O}_X)$ \emph{defined by the sheaf of ideals $\mathscr{I}$};
this is a special case of the general notion of \emph{sub-prescheme}:

\begin{definition}[4.1.3]
\label{I.4.1.3}
A ringed space $(Y,\mathscr{O}_Y)$ is said to be a
\emph{sub-prescheme} of a prescheme $(X,\mathscr{O}_X)$ if:
\begin{enumerate}
  \item[$1^\circ$] $Y$ is a locally closed subspace of $X$;
  \item[$2^\circ$] If $U$ denotes the largest open set of $X$ containing $Y$
  and such that $Y$ is closed in $U$ (in other words, the complement in
  $X$ of the boundary of $Y$ relative to $\overline{Y}$),
  $(Y,\mathscr{O}_Y)$ is a sub-prescheme of
  $(U,\mathscr{O}_X|U)$ defined by a quasi-coherent sheaf of ideals of
  $\mathscr{O}_X|U$.
\end{enumerate}
The sub-prescheme $(Y,\mathscr{O}_Y)$ of $(X,\mathscr{O}_X)$ is said to be
\emph{closed} if $Y$ is closed in $X$ (in which case $U=X$).
\end{definition}

It follows at once from this definition and from
(\hyperref[I.4.1.2]{4.1.2}) that the closed sub-preschemes of $X$ are in
\emph{canonical one-to-one correspondence} with the quasi-coherent sheaves of
ideals $\mathscr{J}$ of $\mathscr{O}_X$, since the fact that two such
sheaves $\mathscr{J}$, $\mathscr{J}'$ have the same (closed) support $Y$ and
are such that the restrictions of $\mathscr{O}_X/\mathscr{J}$ and of
$\mathscr{O}_X/\mathscr{J}'$ to $Y$ are identical implies at once that
$\mathscr{J}'=\mathscr{J}$.

\begin{env}[4.1.4]
\label{I.4.1.4}
Let $(Y,\mathscr{O}_Y)$ be a sub-prescheme of $X$, $U$ the largest open set
of $X$ containing $Y$ and in which $Y$ is closed, and $V$ an open set of $X$
contained in $U$; then $V\cap Y$ is closed in $V$. Moreover, if $Y$ is
defined by the quasi-coherent sheaf $\mathscr{J}$ of ideals of
$\mathscr{O}_X|U$, then $\mathscr{J}|V$ is a quasi-coherent sheaf of ideals of
$\mathscr{O}_X|V$, and it is immediate that the prescheme induced by $Y$ on
$Y\cap V$ is the closed sub-prescheme of $V$ defined by the sheaf of ideals
$\mathscr{J}|V$. Conversely:
\end{env}

\begin{proposition}[4.1.5]
\label{I.4.1.5}
Let $(Y,\mathscr{O}_Y)$ be a ringed space such that $Y$ is a subspace of
$X$ and such that there exists a covering $(V_\alpha)$ of $Y$ by open sets of
$X$ such that for every $\alpha$, $Y\cap V_\alpha$ is closed in $V_\alpha$
and the ringed space
$(Y\cap V_\alpha,\mathscr{O}_Y|(Y\cap V_\alpha))$ is a closed
sub-prescheme of the prescheme induced on $V_\alpha$ by $X$. Then
$(Y,\mathscr{O}_Y)$ is a sub-prescheme of $X$.
\end{proposition}

The hypothesis implies that $Y$ is locally closed in $X$ and that the largest
open set $U$ containing $Y$ and in which $Y$ is closed contains all the
$V_\alpha$; we can therefore reduce to the case where $U=X$ and $Y$ is closed
in $X$. We then define a quasi-coherent sheaf of ideals $\mathscr{J}$ of
$\mathscr{O}_X$ by taking for $\mathscr{J}|V_\alpha$ the sheaf of ideals of
$\mathscr{O}_X|V_\alpha$ which defines the closed sub-prescheme
$(Y\cap V_\alpha,\mathscr{O}_Y|(Y\cap V_\alpha))$ and, for every open set $W$
of $X$ not meeting $Y$, $\mathscr{J}|W=\mathscr{O}_X|W$. One verifies
immediately, by virtue of (\hyperref[I.4.1.3]{4.1.3}) and
(\hyperref[I.4.1.4]{4.1.4}), that there exists one and only one sheaf of
ideals $\mathscr{J}$ satisfying these conditions and that it defines the closed
sub-prescheme $(Y,\mathscr{O}_Y)$.

In particular, the prescheme \emph{induced} by $X$ on an \emph{open set} of
$X$ is a \emph{sub-prescheme} of $X$.

\begin{proposition}[4.1.6]
\label{I.4.1.6}
A sub-prescheme (resp. a closed sub-prescheme) of a
sub\oldpage[I]{121}prescheme (resp. closed sub-prescheme) of $X$
is canonically identified with a sub-prescheme (resp. closed sub-prescheme) of
$X$.
\end{proposition}

Since a locally closed subset of a locally closed subspace of $X$ is
a locally closed subspace of $X$, it is clear
(\hyperref[I.4.1.5]{4.1.5}) that the question is local and that we can
therefore assume $X$ affine; the proposition then follows from the canonical
identification of $A/\mathfrak{J}'$ and
$(A/\mathfrak{J})/(\mathfrak{J}'/\mathfrak{J})$ when
$\mathfrak{J}$, $\mathfrak{J}'$ are two ideals of a ring $A$ such that
$\mathfrak{J}\subset\mathfrak{J}'$.

We shall always make the preceding identification from now on.

\begin{env}[4.1.7]
\label{I.4.1.7}
Let $Y$ be a sub-prescheme of a prescheme $X$, and denote by $\psi$
the canonical injection $Y\to X$ of the \emph{underlying spaces}; we know that
the inverse image $\psi^*(\mathscr{O}_X)$ is the restriction
$\mathscr{O}_X|Y$ (\hyperref[0.3.7.1]{0, 3.7.1}). If, for every $y\in Y$,
we denote by $\omega_y$ the canonical homomorphism
$(\mathscr{O}_X)_y\to(\mathscr{O}_Y)_y$, these homomorphisms are the
restrictions to the stalks of a \emph{surjective} homomorphism $\omega$ of
sheaves of rings $\mathscr{O}_X|Y\to\mathscr{O}_Y$: indeed, it suffices to
verify this locally on $Y$, that is to say, we can assume $X$
affine and the sub-prescheme $Y$ closed; if in this case $\mathscr{I}$ is the
sheaf of ideals in $\mathscr{O}_X$ which defines $Y$, the $\omega_y$ are
nothing but the restrictions to the stalks of the homomorphism
$\mathscr{O}_X|Y\to(\mathscr{O}_X/\mathscr{I})|Y$.

We have thus defined a \emph{monomorphism of ringed spaces}
(\hyperref[0.4.1.1]{0, 4.1.1}) $j=(\psi,\omega^b)$, which is obviously a
morphism $Y\to X$ of preschemes (\hyperref[I.2.2.1]{2.2.1}), and which we
shall call the \emph{canonical injection morphism}.

If $f:X\to Z$ is a morphism, we shall also say that the composite morphism
$Y\xrightarrow{j}X\xrightarrow{f}Z$ is the \emph{restriction} of $f$ to the
sub-prescheme $Y$.
\end{env}

\begin{env}[4.1.8]
\label{I.4.1.8}
In accordance with the general definitions (T, I, 1.1), we shall say that a
morphism of preschemes $f:Z\to X$ is \emph{majorized} by the injection
morphism $j:Y\to X$ of a sub-prescheme $Y$ of $X$ if $f$ factorizes as
$Z\xrightarrow{g}Y\xrightarrow{j}X$, where $g$ is a morphism of preschemes;
$g$ is necessarily \emph{unique}, since $j$ is a monomorphism.
\end{env}

\begin{proposition}[4.1.9]
\label{I.4.1.9}
For a morphism $f:Z\to X$ to be \emph{majorized by an injection
morphism} $j:Y\to X$, it is necessary and sufficient that $f(Z)\subset Y$
and that, for every $z\in Z$, if we set $y=f(z)$, the homomorphism
$(\mathscr{O}_X)_y\to\mathscr{O}_z$ corresponding to $f$ factorize as
$(\mathscr{O}_X)_y\to(\mathscr{O}_Y)_y\to\mathscr{O}_z$ (or, what amounts
to the same, that the kernel of $(\mathscr{O}_X)_y\to\mathscr{O}_z$ contain
that of $(\mathscr{O}_X)_y\to(\mathscr{O}_Y)_y$).
\end{proposition}

The conditions are obviously necessary. To see that they are
sufficient, we can reduce to the case where $Y$ is a \emph{closed}
sub-prescheme of $X$, by replacing $X$ if necessary by an open set $U$ such that
$Y$ is closed in $U$ (\hyperref[I.4.1.3]{4.1.3}); $Y$ is then defined
by a quasi-coherent sheaf $\mathscr{I}$ of ideals of $\mathscr{O}_X$.
Set $f=(\psi,\theta)$, and let $\mathscr{J}$ be the sheaf of ideals of
$\psi^*(\mathscr{O}_X)$ which is the kernel of
$\theta^\sharp:\psi^*(\mathscr{O}_X)\to\mathscr{O}_Z$; taking into account the
properties of the functor $\psi^*$ (\hyperref[0.3.7.2]{0, 3.7.2}),
the hypothesis implies that for every $z\in Z$ we have
$(\psi^*(\mathscr{I}))_z\subset\mathscr{J}_z$, and consequently
$\psi^*(\mathscr{I})\subset\mathscr{J}$. Hence $\theta^\sharp$
factorizes as
\[
  \psi^*(\mathscr{O}_X)\longrightarrow
  \psi^*(\mathscr{O}_X)/\psi^*(\mathscr{I})
  =\psi^*(\mathscr{O}_X/\mathscr{I})
  \xrightarrow{\omega}\mathscr{O}_Z,
\]
the first arrow being the canonical homomorphism. Let $\psi'$ be the continuous
map $Z\to Y$ coinciding with $\psi$; it is clear that we have
${\psi'}^*(\mathscr{O}_Y)=\psi^*(\mathscr{O}_X/\mathscr{I})$; on the other
hand, $\omega$ is obviously a local homomorphism, so
$g=(\psi',\omega^b)$ is a morphism $Z\to Y$
\oldpage[I]{122}
of preschemes (\hyperref[I.2.2.1]{2.2.1}) which, by virtue of the
foregoing, is such that $f=j\circ g$, whence the proposition.

\begin{corollary}[4.1.10]
\label{I.4.1.10}
For an injection morphism $Z\to X$ to be \emph{majorized by the injection
morphism} $Y\to X$, it is necessary and sufficient that $Z$ be a sub-prescheme
of $Y$.
\end{corollary}

We then write $Z\leq Y$, and this relation is obviously an
\emph{order relation} on the set of sub-preschemes of $X$.

\subsection{Immersion morphisms}
\label{subsection:I.4.2}

\begin{definition}[4.2.1]
\label{I.4.2.1}
A morphism $f:Y\to X$ is said to be an \emph{immersion} (resp. a
\emph{closed immersion}, an \emph{open immersion}) if it factorizes as
$Y\xrightarrow{g}Z\xrightarrow{j}X$, where $g$ is an isomorphism, $Z$ a
sub-prescheme of $X$ (resp. a closed sub-prescheme, a sub-prescheme induced
on an open set) and $j$ the injection morphism.
\end{definition}

The sub-prescheme $Z$ and the isomorphism $g$ are then determined
\emph{uniquely}, for if $Z'$ is a second sub-prescheme of $X$, $j'$ the
injection $Z'\to X$ and $g'$ an isomorphism $Y\to Z'$ such that
$j\circ g=j'\circ g'$, we deduce $j'=j\circ g\circ {g'}^{-1}$, whence
$Z'\leq Z$ (\hyperref[I.4.1.10]{4.1.10}), and one shows likewise that
$Z\leq Z'$, so $Z'=Z$, and since $j$ is a monomorphism of preschemes,
$g'=g$.

We say that $f=j\circ g$ is the \emph{canonical factorization} of the
immersion $f$, and the sub-prescheme $Z$ and the isomorphism $g$ are said to be
\emph{associated with $f$}.

It is clear that an immersion is a \emph{monomorphism} of preschemes
(\hyperref[I.4.1.7]{4.1.7}) and \emph{a fortiori} a \emph{radicial}
morphism (\hyperref[I.3.5.4]{3.5.4}).

\begin{proposition}[4.2.2]
\label{I.4.2.2}
\begin{enumerate}
  \item[a)] For a morphism $f=(\psi,\theta):Y\to X$ to be an
  \emph{open immersion}, it is necessary and sufficient that $\psi$ be a
  homeomorphism of $Y$ onto an open subset of $X$, and that for every
  $y\in Y$ the homomorphism
  $\theta_y^\sharp:\mathscr{O}_{\psi(y)}\to\mathscr{O}_y$ be bijective.
  \item[b)] For a morphism $f=(\psi,\theta):Y\to X$ to be an
  \emph{immersion} (resp. a \emph{closed immersion}), it is necessary and
  sufficient that $\psi$ be a homeomorphism of $Y$ onto a locally closed
  (resp. closed) subset of $X$, and that for every $y\in Y$ the homomorphism
  $\theta_y^\sharp:\mathscr{O}_{\psi(y)}\to\mathscr{O}_y$ be surjective.
\end{enumerate}
\end{proposition}

\begin{enumerate}
  \item[a)] The conditions are obviously necessary. Conversely, if they
  are satisfied, it is clear that $\theta^\sharp$ is an isomorphism of
  $\psi^*(\mathscr{O}_X)$ onto $\mathscr{O}_Y$, and
  $\psi^*(\mathscr{O}_X)$ is the sheaf deduced by transport of structure by
  means of $\psi^{-1}$ from $\mathscr{O}_X|\psi(Y)$; whence the
  conclusion.
  \item[b)] The conditions being obviously necessary, let us prove that they
  are sufficient. Consider first the special case where it is assumed that
  $X$ is an affine scheme and that $Z=\psi(Y)$ is \emph{closed} in $X$. We
  know (\hyperref[0.3.4.6]{0, 3.4.6}) that $\psi_*(\mathscr{O}_Y)$ then has
  support $Z$ and that, if we denote by $\mathscr{O}'_Z$ its restriction
  to $Z$, the ringed space $(Z,\mathscr{O}'_Z)$ is deduced from
  $(Y,\mathscr{O}_Y)$ by transport of structure by means of the
  homeomorphism $\psi$, regarded as a map of $Y$ onto $Z$.
  Let us show that $f_*(\mathscr{O}_Y)=\psi_*(\mathscr{O}_Y)$ is a
  \emph{quasi-coherent} $\mathscr{O}_X$-module. Indeed, for every
  $x\notin Z$, $\psi_*(\mathscr{O}_Y)$, restricted to a suitable neighbourhood
  of $x$, is zero. If on the contrary $x\in Z$, we have $x=\psi(y)$ for a
  well-determined $y\in Y$; let $V$ be an affine open neighbourhood of $y$ in
  $Y$; $\psi(V)$ is then open in $Z$, hence is the trace on $Z$ of an open set
  $U$ of $X$, and the restriction to $U$ of $\psi_*(\mathscr{O}_Y)$ is identical
  to the restriction to $U$ of the direct
\oldpage[I]{123}
  image $(\psi_V)_*(\mathscr{O}_Y|V)$, where $\psi_V$ is the restriction of
  $\psi$ to $V$. Now, the restriction to $(V,\mathscr{O}_Y|V)$ of the morphism
  $(\psi,\theta)$ is a morphism of this prescheme into
  $(X,\mathscr{O}_X)$, and is consequently of the form
  $({}^a\varphi,\widetilde{\varphi})$, where $\varphi$ is a homomorphism of
  the ring $A=\Gamma(X,\mathscr{O}_X)$ into the ring
  $\Gamma(V,\mathscr{O}_Y)$ (\hyperref[I.1.7.3]{1.7.3}); we conclude
  that $(\psi_V)_*(\mathscr{O}_Y|V)$ is a quasi-coherent $\mathscr{O}_X$-module
  (\hyperref[I.1.6.3]{1.6.3}), which proves our
  assertion, because of the local character of quasi-coherent sheaves.
  Moreover, the hypothesis that $\psi$ is a homeomorphism implies
  (\hyperref[0.3.4.5]{0, 3.4.5}) that for every $y\in Y$, $\psi_y$ is an
  isomorphism
  $(\psi_*(\mathscr{O}_Y))_{\psi(y)}\to\mathscr{O}_y$; since the diagram
  \phantomsection\label{I.4.2.2.theta-diagram}
  \[
    \xymatrix{
      \mathscr{O}_{\psi(y)}\ar[r]^{\theta_{\psi(y)}}
        \ar[d]_{\psi_y\circ\rho_{\psi(y)}} &
      (\psi_*(\mathscr{O}_Y))_{\psi(y)}\ar[d]^{\psi_y}\\
      (\psi^*(\mathscr{O}_X))_y\ar[r]^{\theta_y^\sharp} &
      \mathscr{O}_y
    }
  \]
  is commutative and the two vertical arrows are isomorphisms
  (\hyperref[0.3.7.2]{0, 3.7.2}), the hypothesis that $\theta_y^\sharp$ is
  surjective implies that so is $\theta_{\psi(y)}$. Since the
  support of $\psi_*(\mathscr{O}_Y)$ is $Z=\psi(Y)$, $\theta$ is a
  \emph{surjective} homomorphism of $\mathscr{O}_X=\widetilde{A}$ into the
  quasi-coherent $\mathscr{O}_X$-module $f_*(\mathscr{O}_Y)$. Consequently there
  exists a unique isomorphism $\omega$ of a quotient sheaf
  $\widetilde{A}/\widetilde{\mathfrak{J}}$ ($\mathfrak{J}$ an ideal of $A$) onto
  $f_*(\mathscr{O}_Y)$ which, composed with the canonical homomorphism
  $\widetilde{A}\to\widetilde{A}/\widetilde{\mathfrak{J}}$, gives $\theta$
  (\hyperref[I.1.3.8]{1.3.8}); if $\mathscr{O}_Z$ denotes the restriction
  of $\widetilde{A}/\widetilde{\mathfrak{J}}$ to $Z$, then
  $(Z,\mathscr{O}_Z)$ is a sub-prescheme of $(X,\mathscr{O}_X)$, and $f$
  factorizes into the canonical injection of this sub-prescheme into $X$ and
  the isomorphism $(\psi_0,\omega_0)$, where $\psi_0$ is $\psi$ regarded as a
  map of $Y$ onto $Z$, and $\omega_0$ the restriction of $\omega$ to
  $\mathscr{O}_Z$.

  Let us pass to the general case. Let $U$ be an affine open set in $X$ such
  that $U\cap\psi(Y)$ is closed in $U$ and non-empty. By restricting $f$ to the
  prescheme induced by $Y$ on the open set $\psi^{-1}(U)$, and regarding it
  as a morphism of this prescheme into the prescheme induced by $X$ on
  $U$, we are reduced to the first case; the restriction of $f$ to
  $\psi^{-1}(U)$ is therefore a closed immersion $\psi^{-1}(U)\to U$,
  factorizing canonically as $j_U\circ g_U$, where $g_U$ is an isomorphism
  of the prescheme $\psi^{-1}(U)$ onto a sub-prescheme $Z_U$ of $U$, and $j_U$
  the canonical injection $Z_U\to U$. Let $V$ be a second affine open set of
  $X$ such that $V\subset U$; since the restriction $Z'_V$ of $Z_U$ to $V$ is
  a sub-prescheme of the prescheme $V$, the restriction of $f$ to $\psi^{-1}(V)$
  factorizes as $j'_V\circ g'_V$, where $j'_V$ is the canonical injection
  $Z'_V\to V$ and $g'_V$ an isomorphism of $\psi^{-1}(V)$ onto $Z'_V$. By
  the uniqueness of the canonical factorization of an immersion
  (\hyperref[I.4.2.1]{4.2.1}), we necessarily have $Z'_V=Z_V$ and
  $g'_V=g_V$. We conclude (\hyperref[I.4.1.5]{4.1.5}) that there is a
  sub-prescheme $Z$ of $X$ whose underlying space is $\psi(Y)$ and whose
  restriction to each $U\cap\psi(Y)$ is $Z_U$; the $g_U$ are then the
  restrictions to the $\psi^{-1}(U)$ of an isomorphism $g:Y\to Z$ such that
  $f=j\circ g$, where $j$ is the canonical injection $Z\to X$.
\end{enumerate}

\begin{corollary}[4.2.3]
\label{I.4.2.3}
Let $X$ be an affine scheme. For a morphism
$f=(\psi,\theta):Y\to X$ to be a closed immersion, it is necessary and
sufficient that $Y$ be an affine scheme and that the homomorphism
$\Gamma(\theta):\Gamma(\mathscr{O}_X)\to\Gamma(\mathscr{O}_Y)$ be surjective.
\end{corollary}

\begin{corollary}[4.2.4]
\label{I.4.2.4}
\begin{enumerate}
  \item[a)] Let $f$ be a morphism $Y\to X$, and $(V_\lambda)$ a covering
  of $f(Y)$ by open sets of $X$. For $f$ to be an immersion (resp. an
  open immersion), it is necessary and sufficient
\oldpage[I]{124}
  that its restriction to each of the induced preschemes $f^{-1}(V_\lambda)$ be
  an immersion (resp. an open immersion) into $V_\lambda$.
  \item[b)] Let $f$ be a morphism $Y\to X$, and $(V_\lambda)$ an open covering
  of $X$. For $f$ to be a closed immersion, it is necessary and sufficient
  that its restriction to each of the induced preschemes $f^{-1}(V_\lambda)$ be
  a closed immersion into $V_\lambda$.
\end{enumerate}
\end{corollary}

Let $f=(\psi,\theta)$; in case a), $\theta_y^\sharp$ is surjective
(resp. bijective) for every $y\in Y$, and in case b) $\theta_y^\sharp$ is
surjective for every $y\in Y$; it therefore suffices to verify that $\psi$, in
case a), is a homeomorphism of $Y$ onto a locally closed (resp.
open) subset of $X$, and in case b), a homeomorphism of $Y$ onto a closed
subset of $X$. Now, $\psi$ is obviously injective and transforms every
neighbourhood of $y$ in $Y$ into a neighbourhood of $\psi(y)$ in $\psi(Y)$ for
every $y\in Y$, by virtue of the hypothesis; in case a), $\psi(Y)\cap V_\lambda$
is locally closed (resp. open) in $V_\lambda$, so $\psi(Y)$ is locally closed
(resp. open) in the union of the $V_\lambda$, and \emph{a fortiori} in $X$;
in case b), $\psi(Y)\cap V_\lambda$ is closed in $V_\lambda$, so
$\psi(Y)$ is closed in $X$ since $X=\bigcup_\lambda V_\lambda$.

\begin{proposition}[4.2.5]
\label{I.4.2.5}
The composite of two immersions (resp. of two open immersions, of two
closed immersions) is an immersion (resp. an open immersion, a
closed immersion).
\end{proposition}

This follows trivially from (\hyperref[I.4.1.6]{4.1.6}).

\subsection{Product of immersions.}
\label{subsection:I.4.3}

\begin{proposition}[4.3.1]
\label{I.4.3.1}
Let $\alpha:X'\to X$, $\beta:Y'\to Y$ be two $S$-morphisms; if $\alpha$ and
$\beta$ are immersions (resp. open immersions, closed
immersions), then $\alpha\times_S\beta$ is an immersion (resp. an open
immersion, a closed immersion). Moreover, if $\alpha$ (resp. $\beta$)
identifies $X'$ (resp. $Y'$) with a sub-prescheme $X''$ (resp. $Y''$) of $X$
(resp. $Y$), then $\alpha\times_S\beta$ identifies the underlying space of
$X'\times_S Y'$ with the subspace $p^{-1}(X'')\cap q^{-1}(Y'')$ of the
underlying space of $X\times_S Y$, where $p$ and $q$ denote the projections of
$X\times_S Y$ to $X$ and $Y$ respectively.
\end{proposition}

Taking into account Definition (\hyperref[I.4.2.1]{4.2.1}), we can restrict
ourselves to the case where $X'$ and $Y'$ are sub-preschemes, and $\alpha$ and
$\beta$ the injection morphisms. The proposition has already been established
for the sub-preschemes induced on open sets (\hyperref[I.3.2.7]{3.2.7});
since every sub-prescheme is a closed sub-prescheme of a prescheme induced on an
open set (\hyperref[I.4.1.3]{4.1.3}), we are reduced to the case where $X'$
and $Y'$ are \emph{closed} sub-preschemes.

Let us first show that we can assume that $S$ is \emph{affine}. Indeed, let
$(S_\lambda)$ be a covering of $S$ consisting of affine open sets; if
$\varphi$ and $\psi$ are the structural morphisms of $X$ and $Y$, let
$X_\lambda=\varphi^{-1}(S_\lambda)$, $Y_\lambda=\psi^{-1}(S_\lambda)$. The
restriction $X'_\lambda$ (resp. $Y'_\lambda$) of $X'$ (resp. $Y'$) to
$X_\lambda\cap X'$ (resp. $Y_\lambda\cap Y'$) is a closed sub-prescheme of
$X_\lambda$ (resp. $Y_\lambda$), the preschemes $X_\lambda$, $Y_\lambda$,
$X'_\lambda$, $Y'_\lambda$ can be regarded as
$S_\lambda$-preschemes, and the products $X_\lambda\times_S Y_\lambda$ and
$X_\lambda\times_{S_\lambda}Y_\lambda$ (resp.
$X'_\lambda\times_S Y'_\lambda$ and
$X'_\lambda\times_{S_\lambda}Y'_\lambda$) are identical
(\hyperref[I.3.2.5]{3.2.5}). If the proposition is true when $S$ is
affine, the restriction of $\alpha\times_S\beta$ to each of the
$X'_\lambda\times_S Y'_\lambda$ will therefore be an immersion
(\hyperref[I.3.2.7]{3.2.7}). Since the product
$X'_\lambda\times_S Y'_\mu$ (resp. $X_\lambda\times_S Y_\mu$) is identified
with $(X'_\lambda\cap X'_\mu)\times_S(Y'_\lambda\cap Y'_\mu)$ (resp.
$(X_\lambda\cap X_\mu)\times_S(Y_\lambda\cap Y_\mu)$)
(\hyperref[I.3.2.6.4]{3.2.6.4}), the restriction of
$\alpha\times_S\beta$
\oldpage[I]{125}
to each of the $X'_\lambda\times_S Y'_\mu$ is again an immersion; consequently
the same holds for $\alpha\times_S\beta$ by virtue of
(\hyperref[I.4.2.4]{4.2.4}).

In the second place, let us prove that we can also assume that $X$ and $Y$
are \emph{affine}. Indeed, let $(U_i)$ (resp. $(V_j)$) be a covering of $X$
(resp. $Y$) by affine open sets, and let $X'_i$ (resp. $Y'_j$) be the
restriction of $X'$ (resp. $Y'$) to $X'\cap U_i$ (resp. $Y'\cap V_j$), which
is a closed sub-prescheme of $U_i$ (resp. $V_j$); $U_i\times_S V_j$ is
identified with the restriction of $X\times_S Y$ to
$p^{-1}(U_i)\cap q^{-1}(V_j)$ (\hyperref[I.3.2.7]{3.2.7}); and likewise, if
$p'$, $q'$ are the projections of $X'\times_S Y'$,
$X'_i\times_S Y'_j$ is identified with the restriction of $X'\times_S Y'$ to
${p'}^{-1}(X'_i)\cap{q'}^{-1}(Y'_j)$. Set
$\gamma=\alpha\times_S\beta$; by definition we have
$p\circ\gamma=\alpha\circ p'$, $q\circ\gamma=\beta\circ q'$; since
$X'_i=\alpha^{-1}(U_i)$, $Y'_j=\beta^{-1}(V_j)$, we also have
${p'}^{-1}(X'_i)=\gamma^{-1}(p^{-1}(U_i))$,
${q'}^{-1}(Y'_j)=\gamma^{-1}(q^{-1}(V_j))$, whence
\[
  {p'}^{-1}(X'_i)\cap{q'}^{-1}(Y'_j)
  =\gamma^{-1}\bigl(p^{-1}(U_i)\cap q^{-1}(V_j)\bigr)
  =\gamma^{-1}(U_i\times_S V_j)
\]
we conclude as in the first part of the argument.

Suppose therefore that $X$, $Y$, $S$ are affine, and let $B$, $C$, $A$ be their
respective rings. Then $B$ and $C$ are $A$-algebras, and $X'$ and $Y'$ are
affine schemes whose rings are quotient algebras $B'$, $C'$ of $B$ and
$C$ respectively. Moreover, we have
$\alpha=({}^a\rho,\widetilde{\rho})$, $\beta=({}^a\sigma,\widetilde{\sigma})$,
where $\rho$ and $\sigma$ are respectively the canonical homomorphisms
$B\to B'$, $C\to C'$ (\hyperref[I.1.7.3]{1.7.3}). This being so, we know
that $X\times_S Y$ (resp. $X'\times_S Y'$) is an affine scheme with ring
$B\otimes_A C$ (resp. $B'\otimes_A C'$), and
$\alpha\times_S\beta=({}^a\tau,\widetilde{\tau})$, where $\tau$ is
the homomorphism $\rho\otimes\sigma$ of $B\otimes_A C$ into
$B'\otimes_A C'$ (\hyperref[I.3.2.2]{3.2.2} and
\hyperref[I.3.2.3]{3.2.3}); since this homomorphism is surjective,
$\alpha\times_S\beta$ is indeed an immersion. Moreover, if $\mathfrak{b}$
(resp. $\mathfrak{c}$) is the kernel of $\rho$ (resp. $\sigma$), the kernel of
$\tau$ is the sum of the ideals generated by $u(\mathfrak{b})$ and $v(\mathfrak{c})$, where $u$ (resp. $v$) is
the homomorphism $b\mapsto b\otimes1$ (resp. $c\mapsto1\otimes c$). Since
$p=({}^a u,\widetilde{u})$ and $q=({}^a v,\widetilde{v})$, this kernel
corresponds, in the prime spectrum of $B\otimes_A C$, to the closed set
$p^{-1}(X')\cap q^{-1}(Y')$ (\hyperref[I.1.2.2.1]{1.2.2.1} and
\hyperref[I.1.1.2]{1.1.2 (iii)}), which completes the proof.

\begin{corollary}[4.3.2]
\label{I.4.3.2}
If $f:X\to Y$ is an immersion (resp. an open immersion, a closed
immersion) and an $S$-morphism, then $f_{(S')}$ is an immersion (resp. an open
immersion, a closed immersion) for every extension $S'\to S$ of the base
prescheme.
\end{corollary}

\subsection{Inverse image of a sub-prescheme.}
\label{subsection:I.4.4}

\begin{proposition}[4.4.1]
\label{I.4.4.1}
Let $f:X\to Y$ be a morphism, $Y'$ a sub-prescheme (resp. a
closed sub-prescheme, a prescheme induced on an open set) of $Y$, and
$j:Y'\to Y$ the injection morphism. Then the projection $p:X\times_Y Y'\to X$
is an immersion (resp. a closed immersion, an open immersion); the
sub-prescheme of $X$ associated with $p$ has $f^{-1}(Y')$ as underlying space;
moreover, if $j'$ is the injection morphism of this sub-prescheme into $X$,
then for a morphism $h:Z\to X$ to be such that $f\circ h:Z\to Y$ is majorized
by $j$, it is necessary and sufficient that $h$ be majorized by $j'$.
\end{proposition}

Since $p=1_X\times_Y j$ (\hyperref[I.3.3.4]{3.3.4}), the first assertion
follows from (\hyperref[I.4.3.1]{4.3.1}); the second is a special case
of (\hyperref[I.3.5.10]{3.5.10}) (where the roles of $X$ and
of $Y'$ must be interchanged). Finally, if we have $f\circ h=j\circ h'$, where
$h'$ is a morphism $Z\to Y'$, it follows from the definition of the product
that we have $h=p\circ u$, where $u$ is a morphism $Z\to X\times_Y Y'$, whence
the last assertion.

We shall say that the sub-prescheme of $X$ thus defined is the \emph{inverse
image} of the sub-prescheme $Y'$ of $Y$ under the morphism $f$, a terminology
which agrees with the one introduced
\oldpage[I]{126}
more generally in (\hyperref[I.3.3.6]{3.3.6}). When we
speak of $f^{-1}(Y')$ as a sub-prescheme of $X$, it will always be
this sub-prescheme that is meant.

When the preschemes $f^{-1}(Y')$ and $X$ are identical, $j'$ is
the identity and every morphism $h:Z\to X$ is therefore majorized by $j'$, so
\emph{the morphism $f:X\to Y$ then factorizes as}
$X\xrightarrow{g}Y'\xrightarrow{j}Y$.

When $y$ is a \emph{closed} point of $Y$ and
$Y'=\operatorname{Spec}(k(y))$ the smallest closed sub-prescheme of $Y$ having
$\{y\}$ as underlying space (\hyperref[I.4.1.9]{4.1.9}), the
closed sub-prescheme $f^{-1}(Y')$ is canonically isomorphic to the
\emph{fibre} $f^{-1}(y)$ defined in (\hyperref[I.3.6.2]{3.6.2}), with
which we shall identify it.

\begin{corollary}[4.4.2]
\label{I.4.4.2}
Let $f:X\to Y$, $g:Y\to Z$ be two morphisms, and $h=g\circ f$ their
composite. For every sub-prescheme $Z'$ of $Z$, the sub-preschemes
$f^{-1}(g^{-1}(Z'))$ and $h^{-1}(Z')$ of $X$ are identical.
\end{corollary}

This follows from the existence of the canonical isomorphism
$X\times_Y(Y\times_Z Z')\xrightarrow{\sim}X\times_Z Z'$
(\hyperref[I.3.3.9.1]{3.3.9.1}).

\begin{corollary}[4.4.3]
\label{I.4.4.3}
Let $X'$, $X''$ be two sub-preschemes of $X$, and $j':X'\to X$,
$j'':X''\to X$ the injection morphisms; then ${j'}^{-1}(X'')$ and
${j''}^{-1}(X')$ are both equal to the greatest lower bound
$\inf(X',X'')$ of $X'$ and $X''$ for the order relation between
sub-preschemes, and canonically isomorphic to $X'\times_X X''$.
\end{corollary}

This follows at once from (\hyperref[I.4.4.1]{4.4.1}) and
(\hyperref[I.4.1.10]{4.1.10}).

\begin{corollary}[4.4.4]
\label{I.4.4.4}
Let $f:X\to Y$ be a morphism, and $Y'$, $Y''$ two sub-preschemes of $Y$;
we have
\[
  f^{-1}(\inf(Y',Y''))=\inf(f^{-1}(Y'),f^{-1}(Y'')).
\]
\end{corollary}

This follows from the existence of the canonical isomorphism between
$(X\times_Y Y')\times_X(X\times_Y Y'')$ and
$X\times_Y(Y'\times_Y Y'')$ (\hyperref[I.3.3.9.1]{3.3.9.1}).

\begin{proposition}[4.4.5]
\label{I.4.4.5}
Let $f:X\to Y$ be a morphism, and $Y'$ a closed sub-prescheme of $Y$ defined
by a quasi-coherent sheaf $\mathscr{K}$ of ideals of $\mathscr{O}_Y$
(\hyperref[I.4.1.3]{4.1.3}); then the closed sub-prescheme $f^{-1}(Y')$ of
$X$ is defined by the quasi-coherent sheaf of ideals
$f^*(\mathscr{K})\mathscr{O}_X$ of $\mathscr{O}_X$.
\end{proposition}

The question is obviously local on $X$ and $Y$; it then suffices to
note that if $A$ is a $B$-algebra and $\mathfrak{K}$ an ideal of $B$,
we have $A\otimes_B(B/\mathfrak{K})=A/\mathfrak{K}A$, and to apply
(\hyperref[I.1.6.9]{1.6.9}).

\begin{corollary}[4.4.6]
\label{I.4.4.6}
Let $X'$ be a closed sub-prescheme of $X$ defined by a quasi-coherent
sheaf of ideals $\mathscr{J}$ of $\mathscr{O}_X$, and $i$ the injection
$X'\to X$; for the restriction $f\circ i$ of $f$ to $X'$ to be majorized
by the injection $j:Y'\to Y$ (in other words, to factorize as $j\circ g$,
where $g$ is a morphism $X'\to Y'$), it is necessary and sufficient that
$f^*(\mathscr{K})\mathscr{O}_X\subset\mathscr{J}$.
\end{corollary}

It suffices to apply Proposition (\hyperref[I.4.4.1]{4.4.1}) to $i$, taking
into account (\hyperref[I.4.4.5]{4.4.5}).

\subsection{Local immersions and local isomorphisms.}
\label{subsection:I.4.5}

\begin{definition}[4.5.1]
\label{I.4.5.1}
Let $f:X\to Y$ be a morphism of preschemes. We say that $f$ is a
\emph{local immersion at a point $x\in X$} if there exist an open
neighbourhood $U$ of $x$ in $X$ and an open neighbourhood $V$ of $f(x)$ in $Y$
such that the restriction of $f$ to the induced prescheme $U$ is a closed
immersion of $U$ into the induced prescheme $V$. We say that $f$ is a local
immersion if $f$ is a local immersion at every point of $X$.
\end{definition}

\begin{definition}[4.5.2]
\label{I.4.5.2}
A morphism $f:X\to Y$ is said to be a local isomorphism at
\oldpage[I]{127}
a point $x\in X$ if there exists an open neighbourhood $U$ of $x$ in $X$ such
that the restriction of $f$ to the induced prescheme $U$ is an open immersion
of $U$ into $Y$. We say that $f$ is a local isomorphism if $f$ is a
local isomorphism at every point of $X$.
\end{definition}

\begin{env}[4.5.3]
\label{I.4.5.3}
An immersion (resp. a closed immersion) $f:X\to Y$ can therefore be
characterized as a local immersion such that $f$ is a homeomorphism
of the underlying space of $X$ onto a subset (resp. a closed subset) of
$Y$. An open immersion $f$ can be characterized as an \emph{injective} local
isomorphism.
\end{env}

\begin{proposition}[4.5.4]
\label{I.4.5.4}
Let $X$ be an irreducible prescheme, and $f:X\to Y$ an injective dominant
morphism. If $f$ is a local immersion, then $f$ is an immersion and $f(X)$
is open in $Y$.
\end{proposition}

Indeed, let $x\in X$, and let $U$ be an open neighbourhood of $x$ and $V$ an
open neighbourhood of $f(x)$ in $Y$ such that the restriction of $f$ to $U$
is a closed immersion into $V$; since $U$ is dense in $X$, $f(U)$
is dense in $Y$ by hypothesis, so $f(U)=V$ and $f$ is a homeomorphism
of $U$ onto $V$; the hypothesis that $f$ is injective implies that
$f^{-1}(V)=U$, whence the proposition at once.

\begin{proposition}[4.5.5]
\label{I.4.5.5}
\begin{enumerate}
  \item[\rm(i)] The composite of two local immersions (resp. of two
    local isomorphisms) is a local immersion (resp. a local
    isomorphism).
  \item[\rm(ii)] Let $f:X\to X'$, $g:Y\to Y'$ be two $S$-morphisms. If
    $f$ and $g$ are local immersions (resp. local isomorphisms),
    so is $f\times_S g$.
  \item[\rm(iii)] If an $S$-morphism $f$ is a local immersion (resp. a
    local isomorphism), so is $f_{(S')}$ for every extension
    $S'\to S$ of the base prescheme.
\end{enumerate}
\end{proposition}

By (\hyperref[I.3.5.1]{3.5.1}), it suffices to prove (i) and (ii).

(i) follows at once from the transitivity of closed (resp.
open) immersions (\hyperref[I.4.2.5]{4.2.5}) and from the fact that if $f$
is a homeomorphism of $X$ onto a closed subset of $Y$, then for every open set
$U\subset X$, $f(U)$ is open in $f(X)$, so there exists an open
subset $V$ of $Y$ such that $f(U)=V\cap f(X)$, and $f(U)$ is consequently
closed in $V$.

To prove (ii), let $z\in X\times_S Y$ and set $z'=(f\times_S g)(z)$; let $p$, $q$ be the projections of $X\times_S Y$, and $p'$,
$q'$ those of $X'\times_S Y'$. By hypothesis there exist open
neighbourhoods $U$, $U'$, $V$, $V'$ of $x=p(z)$, $x'=p'(z')$, $y=q(z)$,
$y'=q'(z')$ respectively, such that the restrictions of $f$ and $g$ to $U$
and $V$ respectively are closed (resp. open) immersions into
$U'$ and $V'$ respectively. Since the underlying space of $U\times_S V$ and
that of $U'\times_S V'$ are identified with the open neighbourhoods
$p^{-1}(U)\cap q^{-1}(V)$ and ${p'}^{-1}(U')\cap{q'}^{-1}(V')$ of $z$ and
$z'$ respectively (\hyperref[I.3.2.7]{3.2.7}), the proposition follows
from (\hyperref[I.4.3.1]{4.3.1}).
